\section{强化学习与 offline RL：数据从哪里来}

这一节建立最少但完整的 RL 背景。读者需要区分 state、action、reward、return、policy 与 dataset support；如果这些概念混在一起，后文很容易把 return-to-go 当成已知未来，或误以为 offline RL 完全消除了探索问题。

\subsection{MDP、policy 与 cumulative return}

标准 Markov decision process（MDP）可以写成 $\mathcal{M}=(\mathcal{S},\mathcal{A},P,r,\gamma)$。在状态 $s_t\in\mathcal{S}$ 下，policy $\pi(a_t\mid s_t)$ 选择动作 $a_t\in\mathcal{A}$；环境按 transition kernel $P(s_{t+1}\mid s_t,a_t)$ 转移，并产生 reward $r_t$。学习目标是最大化期望折扣回报：
\[
J(\pi)=\mathbb{E}_{\tau\sim\pi}\left[\sum_{t=1}^{T}\gamma^{t-1}r_t\right],
\qquad 0<\gamma\le 1.
\]
这里 $\tau=(s_1,a_1,r_1,\ldots,s_T,a_T,r_T)$ 是 trajectory；$\gamma$ 是 discount factor，用于降低远期奖励权重或保证无限时域求和稳定。

\lecturefigure{slide-05-background-reinforcement-learning.jpg}{RL 闭环：agent 依据 state 选择 action，environment 返回 reward 和下一状态。}{Stanford Online 官方视频 04:00。}

\begin{knowledgebox}{术语消化：RL 中最容易混淆的量}
\begin{tabularx}{\textwidth}{L{2.8cm} X X}
\toprule
术语 & 解决的问题 & 与本讲的关系 \\
\midrule
reward $r_t$ & 当前一步得到的标量反馈 & rollout 后用于更新剩余 return-to-go \\
return $G_t$ & 从时刻 $t$ 起的累计未来 reward & 训练数据中可由完整轨迹回算 \\
policy $\pi$ & 给定信息时如何选择 action & Decision Transformer 通过 action prediction 隐式实现 policy \\
value / critic & 估计状态或动作的长期价值 & 原方法不依赖 critic 训练 actor，但后文展示可改为 critic target \\
credit assignment & 判断早期动作应为远期结果承担多少责任 & Key-to-Door 用跨阶段稀疏 reward 检验 \\
\bottomrule
\end{tabularx}
\end{knowledgebox}

\subsection{Offline RL 的固定数据集}

在线 RL 可以在训练期间持续与环境交互；offline RL（离线强化学习）只使用预先收集的固定数据集
\[
\mathcal{D}=\{\tau^{(i)}\}_{i=1}^{N},
\]
其中轨迹可能来自随机 policy、训练到一半的 policy、专家 policy 或多个 policy 的混合。固定数据降低了真实机器人、医疗或工业系统中的交互成本与风险，却带来更尖锐的问题：模型可能在 rollout 中访问数据集几乎没有覆盖的 state-action 区域。

\lecturefigure{slide-06-background-offline-rl.jpg}{Offline RL 只使用预先记录的 logged interactions，不进行新的 trial-and-error。}{Stanford Online 官方视频 05:20。}

\begin{importantbox}{Dataset support（数据支持）}
若数据集在某个状态 $s$ 附近几乎没有动作 $a$ 的样本，就很难可靠估计执行该动作后的结果。模型可以输出语法上合理的 action，但不能凭函数逼近自动创造反事实证据。Offline RL 的核心不是“有一个大数据集”就足够，而是数据是否覆盖部署时需要的行为与状态分布。
\end{importantbox}

\begin{warningbox}{Distribution shift 的两层含义}
第一层是训练数据由 behavior policy $\mu$ 产生，而部署使用 learned policy $\pi$；第二层是一次小动作误差会把 agent 推入新状态，后续预测继续在更陌生的输入上运行。Decision Transformer 避免了显式 Bellman backup，却没有让这两层分布偏移自动消失。
\end{warningbox}

\subsection{为什么讲者说 RL “不扩展”}

2021 年前后的大型 language model 已达到百亿到千亿参数，而常见 RL policy 往往小得多、训练波动也更明显。slide 把差距归因于 RL objective 的不稳定性，这一判断应理解为课程动机：bootstrap target、on-policy distribution change、exploration 和 reward scale 都会增加优化难度；它不是说参数量本身就能衡量 policy 智能。

\lecturefigure{slide-07-motivating-challenge-rl-does-not-scale.jpg}{课程动机：当时的大语言模型远大于常见 RL policy，且训练更稳定。}{Stanford Online 官方视频 07:00。}

\begin{knowledgebox}{读图：先比较训练接口，再比较参数量}
左侧文字比较了 model scale，右侧曲线比较 training stability。真正支持 Decision Transformer 的论点是：如果 action prediction 可以使用标准 MSE 或 cross-entropy，便能复用监督学习的优化与正则化工具。图中并没有证明“模型越大，控制越好”，也没有控制环境、数据量、action space 和评价预算，因此不能把参数差距解释成单一因果关系。
\end{knowledgebox}

\teachervoice{老师强调，offline benchmark 中的数据通常来自在线 RL agent 的 replay buffer，例如训练到中等水平时保存的交互。算法接口可以对数据来源保持 agnostic，但实验结论不能忽略 collection policy；数据的质量和覆盖会直接影响模型能学到什么。}

\subsection{本章小结}

RL 的目标是最大化真实环境中的累计回报；offline RL 则在固定 logged trajectories 上学习 policy。固定数据消除了训练期的新交互，却放大 dataset support 与 distribution shift。Decision Transformer 接下来改变的是训练表述，不是这些统计约束。

